\subsection{Completion of the proof and the finite computations}
\label{subsec:completion}
\begin{proof}[Proof of Theorem~\ref{thm:computer}]
Lemma~\ref{lem:full-response} gives $v\in C^1((1,m))$ and
$v'(\lambda)=-D_\lambda/Z$, with $Z>0$. Fix $1<\lambda<m$.
If $\alpha=1$, Lemma~\ref{lem:calibration} gives $D_\lambda>0$.
If $0<\alpha<1$, the moments of $U$ satisfy
\eqref{eq:hard-calibration}. Proposition~\ref{prop:entire-low-region}
covers $1<\lambda<2/\sqrt3$,
Section~\ref{sec:near-region} covers $2/\sqrt3\le\lambda\le707/500$,
and Section~\ref{sec:middle-high} covers
$707/500\le\lambda<(3+\sqrt5)/2$.
Each proof evaluates $D_\lambda$ under the escape law $\pi_\lambda$.
Together they give $D_\lambda>0$, including at the shared endpoints
$2/\sqrt3$, $707/500$ and $39/20$.

For $(3+\sqrt5)/2\le\lambda<m$, Theorem~\ref{thm:high} gives
$\Dplus v(\lambda)<0$. Differentiability, already proved by
Lemma~\ref{lem:full-response}, identifies this Dini derivative with
$v'(\lambda)$. Thus $v'(\lambda)<0$ at every point of $(1,m)$.
For any $1<\lambda_1<\lambda_2<m$, the mean value theorem on
$[\lambda_1,\lambda_2]$ now gives $v(\lambda_2)<v(\lambda_1)$.
Lemma~\ref{lem:continuity-one} gives continuity at $\lambda=1$.
Together with Theorem~\ref{thm:low}, this proves strict decrease
throughout $(0,m)$.
\end{proof}

The accompanying computations~\cite{Supplement} are organized by the
three regional proofs, with a reading guide that maps each finite
arithmetic lemma to its executable verification and exact data.
The programs reconstruct the rational expressions stated here, check
their polynomial identities and coefficient inequalities using exact
integer or rational arithmetic, and provide the resulting coefficient
arrays. Rational interval decompositions are recorded explicitly when
needed. The programs do not infer signs from floating-point samples.

The preceding proofs show that the denominators are positive on the stated domains and that the coefficient signs imply the required inequalities there. They also bound the infinite tails, verify that the parameters of the escape law lie in those domains, and deduce $D_\lambda>0$ from the resulting bounds.
The supplementary historical source records document the provenance of
the arguments. Their file-identity and dependency checks are separate
from the exact arithmetic verification used for this manuscript.
